% ============================================================
%  Hans Brøchner, "Unconscious and Conscious Human Mental Life"
%  [Det ubevidste og det bevidste menneskelige Sjæleliv]
%  Nyt dansk Maanedsskrift, ed. Vilhelm Møller, vol. 3 (October 1871 --
%  March 1872), printed pp. 42--57; signed "H. Bröchner."
%  TRANSLATION (English).
%  Translated from transcription.tex (Danish) in this directory, whose
%  header records how its text was established (per-batch word check and
%  a full second reading at the image).  Method: ../../../TRANSLATION-
%  PLAYBOOK.md.  No other edition or translation exists or was consulted.
%
%  ---- REGISTER ----
%  Scholarly, moderately literal, readable English; American spelling.
%  Long periods are broken only where English will not carry them, never
%  across a page marker.
%
%  ---- TERMINOLOGY ----
%  Aligned with this repo's translation of Om det Religiøse (1869):
%    Sjæleliv = mental life; sjælelig = psychic (sjæleligt Liv = psychic
%      life; sjælelig-aandelig = psychic-spiritual); Sjæl = soul;
%    det ubevidste / det bevidste = the unconscious / the conscious;
%    Bevidsthedsliv = life of consciousness; Aand = spirit; Væsen =
%      essence; Tilværelse = existence; det Universelle / det Individuelle
%      = the universal / the individual; det Enkelte = the particular;
%    Forestilling / Almeenforestilling = representation / general
%      representation; Anskuelse = intuition; Erkjendelse = cognition;
%    Vexelvirkning / Vexelforhold = interaction / reciprocal relation;
%    Inderlighed = inwardness; Inderliggjørelse = inwardization;
%    Livsfyldes Organisation = organization of the fullness of life.
%  Names as printed (Leibnitz).
%
%  ---- CONVENTIONS ----
%  - Page markers \opage{N} and "% --- p. N ---" at the same joints as in
%    transcription.tex.
%  - The same paragraphs, numbered sections (\secnum), \emph spans
%    (letterspacing in the original), rules and signature as the Danish.
%    Where the Danish spaces only part of a compound, the emphasis falls
%    on the corresponding part of the English (% TR: at the site).
%  - Quotation marks „...“ -> ``...''.  The signature "H. Bröchner." is
%    kept as printed.
%
%  ---- DEFECTS: DANISH AS PRINTED, ENGLISH CORRECT ----
%  p. 53 "Bevisthedsfordobling" rendered as "doubling of consciousness";
%  p. 57 the gap before "Døden" (a lost sort) rendered "[in] death".
%
%  ---- HOW THE TRANSLATION WAS CHECKED (2026-09-27) ----
%  (1) Two translators (pp. 42--49, 50--57), each re-reading its English
%  against the Danish sentence by sentence.  (2) Mechanical audit
%  (tr_audit.py): per printed page the same \opage joints, \emph spans,
%  section numbers, centre blocks and paragraph breaks as
%  transcription.tex, and an English/Danish word ratio within 0.95--1.60:
%  16 pages, 0 flagged (ratios 1.08--1.27).  (3) Terminology harmonized
%  across the two batches (sjælelig) in the caller.  (4) Independent
%  review: two readers who had not translated the pages read every
%  sentence of the Danish against the English.  No omission, addition or
%  error of sense found; 5 corrections applied -- p. 42 "Deeltheden"
%  (dividedness, not "partiality"); p. 45 "dets Heelhed" (its entirety);
%  p. 48 "i det enkelte Øieblik" (at any given moment); p. 53 "ved sig
%  selv" (through itself, i.e. self-subsistent -- not "by itself");
%  p. 54 "Aands Virken" (working, as elsewhere for Virken).
% ============================================================
\documentclass[12pt,a4paper,openany]{book}
\usepackage[utf8]{inputenc}
\usepackage[T1]{fontenc}
\usepackage{libertinus}
\usepackage{libertinust1math}
\usepackage{textalpha}   % Greek: Διαψάλματα (p. 348), μεσότης (p. 361)
\usepackage{graphicx}
\usepackage[english]{babel}
\usepackage[protrusion=true,expansion=true]{microtype}
\usepackage{setspace}
\usepackage[margin=1.2in]{geometry}
\usepackage{enumitem}
\usepackage{fancyhdr}
\setlength{\headheight}{28pt}
\addtolength{\topmargin}{-13.5pt}

% Footnotes as the printing sets them: *) (none occur in this article).
\renewcommand{\thefootnote}{*)}

% The article's numbered sections: the number alone, centred on its line,
% exactly as in transcription.tex.
\newcommand{\secnum}[1]{\par\begin{center}#1\end{center}\par\nobreak}

\usepackage{hyperref}
\hypersetup{
  pdftitle={Unconscious and Conscious Human Mental Life},
  pdfauthor={Hans Brøchner},
  hidelinks
}

\pagestyle{fancy}
\fancyhf{}
\fancyhead[LE,RO]{\thepage}
\fancyhead[RE]{\textit{H. Brøchner: Unconscious and Conscious Human Mental Life}}
\fancyhead[LO]{\textit{Nyt dansk Maanedsskrift, vol.\ 3}}

\setstretch{1.3}

\title{\textbf{Unconscious and Conscious Human Mental Life}\\[8pt]
  \large [Det ubevidste og det bevidste menneskelige Sjæleliv]}
\author{Hans Brøchner}
\date{\textit{Nyt dansk Maanedsskrift}, ed.\ Vilhelm Møller, vol.~3
  (October 1871--March 1872), pp.~42--57\\[10pt]
  \small Translated from the Danish}

% ---- original-page markers ----
% \opage{N} marks where printed page N begins, at the same joint as in
% transcription.tex (mid-sentence where the page turns mid-sentence).
\usepackage{marginnote}
\newcommand{\opage}[1]{\marginnote{\footnotesize\textit{[#1]}}}
\newcommand{\apage}[1]{\marginnote{\footnotesize\textit{[#1]}}}

\begin{document}
\maketitle
\thispagestyle{empty}
\clearpage

% --- p. 42 ---
\opage{42}\begin{center}\textbf{\large Unconscious and Conscious Human Mental Life.}\end{center}

\secnum{1.}

Among all the problems with which the psychology of our
day is occupied with renewed interest, there is none that has a
wider compass and a greater reach than the problem of the relation
between our conscious and our unconscious mental life.
The series of phenomena it touches reaches from the individual's
first coming into being to its cessation, and the forms of life
whose relation it demands be cleared up point back to the
fundamental oppositions of existence itself and to the relation between them.
What constantly, with irresistible power, leads cognition back
to this problem is our spirit's inescapable need to find unity
in ourselves and in existence, and not to come to rest in dividedness
as the final word. It is this need that twice in one century, and in
wholly opposite directions, has led to one-sided formulas for the
unity of man and of existence: to the unity formula of materialism,
which lets the spiritual in existence vanish into moving matter,
and lets the conscious life in us become a result of the combination
of the elements of matter; --- and to the unity formula of spiritualism,
which lets nature vanish into a spirit without nature, and lets the
life of thought in us tear itself loose from its organic conditionedness.
It is this need that, once again, after the impracticability of the
one-sided formulas has become evident, leads to
% --- p. 43 ---
\opage{43}a striving for a unity that does not misjudge the distinctive character
of the oppositions, but would see that character proceed from, and be
joined together again in, a higher, richer uniting principle. ---

It was at the beginning of modern philosophy that the sharp,
dualistic distinction between what the psychology of our time
designates as unconscious and conscious \emph{mental life},
and strives to lead back to unity, was established through a
consequence drawn not from a comprehensive, unbiased experience
but from metaphysical presuppositions. Spirit and matter,
thought and extension, were determined as opposites that did not
meet in common predicates, but only excluded one another.
In man, to be sure, one had been compelled by the force of the facts
to posit spirit and matter, soul and body, as joined and as acting
upon one another; but the connection was expressly determined as
incomprehensible, and only as a composition, not as a real coming
together into unity. Everything that belongs to the soul was to
fall under the determination of thought and consciousness; everything
that belongs to the body, under that of extension and motion, from
which thought and consciousness were excluded. The body was therefore,
even if the soul was to have a ``seat'' in it --- but as a foreign
guest and by an incomprehensible miracle --- conceived only as a
machine, as a ``divine automaton'', and everything that takes place
in it was to happen according to the laws of mechanism. In the animal,
where no soul was recognized, because the animal was denied thought
and consciousness, this relation was seen to stand out most clearly.
The interaction that was posited in man between soul and body, even
if only through an inconsistency, and which could nevertheless be
realized only in the forms of mechanism, here fell away, and every
impulse to an activity here issued from something bodily. One then
stubbornly refused to conceive even those phenomena within the animal
world that most of all recalled those of human consciousness as
expressions of a psychic activity. Thus, when the mistreated animal
expressed its pain by its cry, one boldly denied, trusting in the
unshakable firmness of one's
% --- p. 44 ---
\opage{44}presuppositions, that the cry was a sign of a sensation.
The animal's cry, it was said, was like the tone of a string,
a purely mechanical effect of the blow that struck it.

From this dualistic theory, which sets an insurmountable boundary
between human life and the lower forms of life, and which would unite
into unity, in the human being, what is absolutely opposed, and thus
irreconcilable, and would let what is assigned to absolutely different
spheres --- which must therefore lack any possibility of contact ---
enter into a relation of interaction, cognition of the absurd
consequences and of the insoluble difficulties to which it led was
nevertheless bound soon to lead away; and a freer, more unbiased view
of the phenomena of reality and a sharper testing of the concepts of
the opposites were bound to lead to a widening of the theory's
determinations of spirit and nature, and thereby to pave the way for
the cognition of the unity that joins them together in the whole of
existence and in the forms of human life.

It was from spirit's consciousness of itself as thinking that that
dualistic theory had taken its point of departure when it sought a
secure foothold for cognition; and as spirit, in order to find its
own essential determination, distinguished itself from everything
from which it could abstract, the bodily was separated from it as
wholly foreign to it and as that which had no determination in
common with it. Therefore thought, and what was made one with it,
namely consciousness, was kept apart by an insurmountable chasm from
spirit's opposite; psychic life received only one form, the life of
consciousness, and between this life of consciousness and what took
place in the unconscious no unity could be thought. And yet
consciousness points of necessity to an essential connection with the
unconscious. In its concept there lies an essential relation to that
which is other than the I, in which other the material too is
comprised; but an essential relation, as a necessary relation, includes
the essential connection; only by that connection is the necessity of
the relation explained. And the life of consciousness has a becoming,
and in becoming it proceeds from an unconscious
% --- p. 45 ---
\opage{45}living that is an abiding condition of the reality of consciousness.
By the very concept that was meant to keep the spiritual apart from
the bodily, there is thus a pointing to a unity in the human being of
spirit and nature, of the conscious and the unconscious; and there is
thereby a pointing to a fuller and richer determination of the
opposites than the one that could be given them as abstractly separated.
In order that the unity can be realized for thought, the essence of
nature must not be determined by what expresses the real (the bodily)
only, as it were, at its greatest distance from spirit, by passive
extension; its concept must be formed not from its poorest forms but
from the highest phenomena, in which it expresses its complete essence
most richly and most clearly: from the organic, the living, in which
the ideality of nature comes to light. And the concept of spirit must
not be determined solely from the phenomena in which the connection
with the natural ground has, as it were, become invisible, but from the
whole circle of its forms, and therefore also from those that show the
spiritual in its temporal development as borne and conditioned by a
real organic ground, the conscious life of spirit in unity with an
unconscious psychic life, through which its connection with nature
becomes manifest. From both points of departure, from nature and from
spirit, one will then arrive at the essential unity of the opposites,
and only by the recognition of this essential unity will the facts
be able to be understood in their completeness.

When the opposites in the human being are seen in unity, the concept
of \emph{life} becomes that which binds them together; it is at once
that which expresses the essence of nature and that which forms the
ground of the essence of spirit. We then see life as comprising the
two forms, the organic and the psychic in the narrower sense, unconscious
and conscious life; and as the human being is seen as a unity and the
unity is determined from the standpoint of the highest form, the
term \emph{psychic} life, in the wider sense, is extended to life
in \emph{its entirety}, and \emph{the soul} is seen as the principle of
\emph{all} life-activity in the human. The soul is then one with the
% --- p. 46 ---
\opage{46}life-principle, and this in turn is the ideal principle of unity by
which the organism exists as a totality, as a unity that asserts itself
in its relation to the external and to its organs and functions in
their multiplicity. We do not, then, conceive the soul as something
foreign that is, as it were, an inhabitant of the body, or as something
immaterial that in an incomprehensible way governs the elements of
matter in the body; but we conceive it as the real organism's
\emph{own} ideality, and we do not posit its ideality as synonymous with
an absolute opposition to the material, but as the form of unity's
being \emph{as} unity in the real manifold, which, precisely by being
able to become a member of a totality, of a living whole that is
present \emph{as such} in its \emph{single} members, shows that it
bears within itself the possibility of ideality.

But when, in the concept ``psychic life'', we include both fundamental
forms of life as in a higher unity, it is implied that our conscious
and our unconscious life must stand in a thoroughgoing reciprocal
relation to one another throughout the whole series of the phenomena
of individual mental life; and this thoroughgoing reciprocal relation,
through which the essential unity everywhere asserts itself, we shall
now try to demonstrate in brief strokes, leaving the sphere of abstract
conceptual determinations in order to pass over to that of the facts.

\secnum{2.}

As the first activity of the life-principle that, as the principle of
both forms of life, was determined as a \emph{psychic} principle, we find
the \emph{organizing} activity. As producing an organic totality, not a
heaping together of atoms, this activity shows itself to be an ideal,
psychic activity, but one that is wholly in the form of the unconscious.
It is the presupposition of conscious mental life, in that it develops
the organs by which the reality of the life of consciousness is
conditioned. But it is not merely the temporal presupposition of the
life of consciousness; it is also, after that life has come to reality,
the abiding condition
% --- p. 47 ---
\opage{47}of it, in that it constantly reproduces the organs of the life of
consciousness.

To this organ-forming, purely unconscious psychic activity there are
then joined, by imperceptible, continuous transitions, the most
elementary forms of the life of consciousness in the fetus and in the
independent individual at its first stage of life. We here take the
concept of the life of consciousness in its widest meaning, in which it
comprises not only the stage at which the I clearly distinguishes itself
as I from what is not the I, but also the stage at which consciousness
still only dawns beneath sensation. From the unconscious reflex movements,
through the dimmest forms of organic sensations and sense-sensations,
with their mechanical pressing forward and pressing back, and on to the
conscious sensations, and again from these to the more and more unfolded
functions of consciousness, representation and thought, we have a series
in which the transition takes place through the infinitely small, and
where the boundary of qualitative differences cannot be drawn by way of
experience. At every point in the advancing series there is a relatively
unconscious and a relatively conscious: the unconscious beginning strives
toward the conscious, and the conscious keeps the unconscious within
itself as a remainder or as an element.

How the life of consciousness that develops in this way out of the
unconscious preserves its indissoluble connection with it can be made
clear from many different standpoints. Thus we see that not merely by
chance, but normally and periodically, \emph{an alternation} takes place
between a state in which our consciousness is awake and a state in which
our psychic life is sunk down into the unconscious. In dreamless sleep,
the normal form of sleep, this periodic return of our total life of
consciousness into the unconscious takes place in such a way that the
return shows itself to be unconditionally necessary, overwhelming every
resistance of the will. And the very health of the life of consciousness,
its composure and its clarity, is conditioned by this periodic immersion
in the strengthening bath of the unconscious.

% --- p. 48 ---
\opage{48}We see the same reciprocal relation in the phenomenon
\emph{that every determinate content of consciousness goes back into
the unconscious and comes forth from it again}. What fills
consciousness at any given moment cannot abidingly be its content;
every attempt to hold the particular fast in consciousness beyond a
certain limit of time will inevitably fail; it is, as it were, forcibly
pushed away from there. And that it must be so lies in the nature of
the matter: the particular, precisely as particular, is incongruent
with the generality of consciousness, which comprises everything
particular; consciousness therefore cannot be pressed into its narrow
limits. But what is drawn away from the focal point of consciousness,
and soon moved entirely out of its field of view, is not lost to
consciousness. It goes back into the unconscious depth and is kept
there as the property of spirit. And it is not merely kept, so that
under certain conditions it can again be lifted up to the light and
remembered, but it is drawn into a process that takes place in the self
beneath the level of consciousness: it enters into combinations with the
rest of what is preserved, is fused together with it, and is thereby
nuanced. Work goes on uninterruptedly in secret, and a development can
take place that surprises us when its results come forth before
consciousness, as when during sleep a thought has ripened, a decision
become clear, or when during a longer period of spiritual rest a
process of spiritual assimilation has taken place, the scattered and
foreign material gathered and appropriated. This activity in the
unconscious depth of the soul constantly sends fine effects up even
to the highest and freest activities of thought and will, and the
conscious spirit has no power to interrupt the work and the development
that go on ceaselessly in the seemingly motionless unconscious inwardness.

From a new standpoint we see the essential connection between the two
forms of life when we observe \emph{the mobility of the boundary between
them}. Thus, for example, a movement, an action, which, when it is
unfamiliar to
% --- p. 49 ---
\opage{49}the individual, is performed only by fixing attention on the
particular parts of the organs that contribute to its performance, comes
through habit to be performed unconsciously, in a mechanical way (for
example, the handling of an instrument). Only the general image of the
result, and even that ever more obscured, stands before consciousness:
all the steps through which the result is brought about become
unconscious to us. Within certain limits this last holds of every
voluntary, and as such conscious, movement. The final form of the
movement stands before consciousness as an image of movement; but the
intermediate links by which the movement is realized, the nerve actions
and muscle contractions that are its conditions, are unconscious to us.
Conversely, what in the normal state of our organism is withdrawn from
consciousness can, under abnormal conditions, become conscious to us:
the boundary, that is, is moved in the opposite direction.

If we analyze \emph{the single act of consciousness}, we again find
the conscious and the unconscious going together in unity, in that the
act of consciousness itself is a unity of unconscious elements. If we
choose, for example, the single conscious sensing, this is easily
shown. The single sensation of a tone, for example, which for our
immediate apprehension stands as something purely single, has as its
presuppositions a physical and an organic process in which the elements
are easy to point out. In the physical stimulus that conditions the
determinate sensation of a tone there is a number of waves corresponding
to the distinctive character of the tone, and in each of these waves,
as it advances, countless degrees of condensation. And in the organic
process that is called forth by that stimulus there must be something
corresponding to the stimulus's elements and differences of degree.
But as the organic process is converted into the psychic one, into the
conscious sense-\emph{sensation}, the manifold is gathered up into a
% TR: p. 49 the Danish spaces only „fornemmelse“ in „Sandse\emph{fornemmelse}“; the English carries the emphasis on "sensation" alone.
unity in which the elements do not become conscious as such, but make
themselves felt only through the qualitative determinateness of the
sensation. What holds of the single sense holds of all, insofar as
processes of motion are the physical and organic conditions of the
various sensations. In higher
% --- p. 50 ---
\opage{50}forms of the life of consciousness we again meet the same relation.
When we form a general representation, e.g.\ of a tree or of an
animal, we form it involuntarily, by a mechanical process in which
the individual sense-images become the elements, which produce a
pronounced image where they, as it were, cover one another, but are
blurred where they differ. But this process, which is the
presupposition of the conscious bringing-out of concepts, remains,
together with its elements, in the unconscious; only the result comes
before consciousness. It is this relation between the individual act
of consciousness and its unconscious elements to which \emph{Leibnitz}
first drew attention. In the roar of the sea, which we perceive as a
single sound, he pointed to the sound of the countless individual
waves, which, without becoming conscious to us, gather into a unity
for our ear; in our perception of the color green, which arises from
the mixing of finely ground blue and yellow, he pointed to the
infinitely small unconscious effects of the individual color-elements
in the mixture. In general, he pointed to the infinitely small as the
unconscious element in every determinate conscious magnitude.

This very relation between the element and the summed unity of the
elements points in turn to a new standpoint for the relation between
the unconscious and the conscious. What remains unconscious for us
owing to its quantitative determinateness can be raised to
consciousness in a twofold way. The physical stimulus that calls forth
a movement in the organism --- which we cannot conceive as separate
from a psychic act: an element of sensation --- but not a movement
strong enough to call forth a \emph{conscious} sensation, can by
\emph{an increase in strength} release one. The same can, within
certain limits, be reached by another route: through a directing of
\emph{attention}. We concentrate ourselves consciously in the
particular direction, and the impression that was before too weak to
be noticed, a fleeting auditory impression, for example, is now
perceived clearly and consciously.

% --- p. 51 ---
\opage{51}Finally, we see the connection between the unconscious and
the conscious mental life from the standpoint that a \emph{constant
reciprocal influence} takes place, through which the essential unity
reveals itself. And through this we see a \emph{circuit} brought
about. What is, as it were, deposited from the sphere of conscious
life into that of the unconscious and determines the latter, partly
with visible effects in the organic, works back from there in turn
upon the life of consciousness through the unconscious activity of
the undivided, concentrated psychic-spiritual unity. From this
concentrated unity, which takes up the results of the life of
consciousness, the general direction of the activities of spirit is
unconsciously determined, and from it there issue, in an involuntary
manner, impulses to peculiar combinations of representations
(associations of ideas) and supplements and correctives for cognition
and will. And from that unconscious, to which the moments of the life
of consciousness, as they step out of the sphere of consciousness,
make their contribution, and in which they are worked together, there
takes place not only, under these forms, an incessant
\emph{working-back} upon the conscious; but already \emph{as the
presupposition and foundation of the life of consciousness} the
unconscious life works from the very first upon the conscious, through
the moods and feelings, desires and promptings of drive grounded in
its individual peculiarity, and in a special way through the
connection into which it brings conscious life with the currents of
universal life, with which it stands in closer and more intimate
union. On the continued assimilation through the process of the
unconscious depends the development of conscious life; on the normal
movement of this process, the health and undisturbed harmony of the
life of consciousness.

Thus we are everywhere directed to the unity of the two forms of life,
to their inner connection of essence. At every stage of development
they point to each other: the unconscious to the conscious as its
goal, its completion and its unity; the conscious to the unconscious
as its presupposition, its abiding condition, its element. Different
in scope --- since the life of consciousness rests on the broader base
of unconscious life, so
% --- p. 52 ---
\opage{52}that our I, although it designates our \emph{whole} being in
its real and ideal aspect, nevertheless expresses our total individual
living only in the form of relative abstraction --- they are one in
essence, forms of the activity of the same psychic principle.

\secnum{3.}

When in the foregoing we considered the relation of the two forms of
life, we saw it solely from the standpoint of individual life and as
pointing back to the relation between the side of nature and the side
of spirit in us. But the relation will appear in a new and clearer
light when the standpoint is widened from the individual to the
universal active in the individual, and when we see the reciprocal
relation between the conscious and the unconscious as conditioned by
the relation of the individual \emph{spirit} to a more universal
\emph{spiritual} power.

We already arrive at this relation by going back to the concept of the
concentrated unity of the self, which makes itself felt in an
immediate way in the individual's conscious life. For in this
concentrated unity we can distinguish between the \emph{acquired},
which, as assimilated, has entered into the unity in the form of
immediate being, and the \emph{original}, which is given as an
endowment of essence and from the beginning works unconsciously; and
this originally unconscious working we now seek precisely in the
\emph{highest spiritual} functions: in the working of the \emph{a
priori} and in \emph{genius}, since in the one as in the other we see
something universal making itself felt through what belongs most
inwardly to individuality.

What we, with a name accepted in philosophy, call the \emph{a priori}
designates the original forms in which the peculiar character of our
spiritual essence gives itself expression throughout all its activity
in the various directions of the life of spirit, and gives itself
expression in such a way that these forms show themselves as laid down,
prior to the influence, in the essence that is determined to activity
by the influence. In every domain of the human spirit where a peculiar
activity of consciousness appears, a peculiarly pronounced a priori
% --- p. 53 ---
\opage{53}must make itself felt: in our intuiting, in our thinking, in
our ethical acting and our artistic producing. But this a priori,
which expresses the fundamental peculiarity of our essence, points
back --- since this essence is not through itself but is grounded in
something higher --- to that higher: the fundamental source of
existence; and it becomes --- since what is the ground of our essence
must be of one essence with what is grounded --- an expression (if an
abstract expression) of the essence of the principle of existence
itself. This a priori, which is a determination of essence in us,
comes to be for consciousness --- precisely as a determination of
essence in \emph{spirit}, which by virtue of its freedom is
\emph{itself} to \emph{acquire} what is originally \emph{given} to it
--- only through an activity of spirit: it must come to be for
consciousness through our own energy, merely incited by experience.
But before it comes to consciousness, it works normatively and
determiningly in every act of consciousness; by its hidden working it
% TR: Danish as printed „Bevisthedsfordobling“ (printer's error for „Bevidsthedsfordobling“); rendered in its evident sense, "doubling of consciousness"
prepares the possibility of the doubling of consciousness by which it
can itself become an object for consciousness, and in the very act by
which it comes forth before consciousness it is itself, as normative
for \emph{every} act of spirit, unconsciously determining. It is the
law of our essence that determines the act by which the basic forms of
the active essence are first to be raised to consciousness. But these
forms of the essence are at the same time forms of something higher,
expressions of the essence of the universal principle, and thereby
valid not only for the individual self but for humanity and for
existence.

What holds of the a priori in general holds by the same token of the
pregnant form of it that comes forth in \emph{genius}. In genius there
is an original possibility in the form of the unconscious, which works
irresistibly toward a revelation in consciousness; but even when the
genius becomes most clearly conscious of itself, something unconscious
is nonetheless preserved in the working of genius. The unity of
necessity and freedom, of the working of nature and the energy of
spirit, that characterizes the genius while it instinctively seeks
revelation, characterizes it still in its highest development of
consciousness. Its highest working is in the form of immediacy, an
unconscious creating,
% --- p. 54 ---
\opage{54}which only when it has given its work reality radiates
through that work the light in which it itself becomes partly
transparent. And in the producing of genius --- in which the individual
stamps its spiritual peculiarity in the most intensive way --- the
individual feels itself to be the organ of a higher power by which it
is pervaded. The deed of genius has the validity of the idea; but
thereby it points back to the ideal power from which the gift of
genius has its source.

By what is most original and deepest in the working of the individual
spirit we are thus led back to \emph{the working of the universal in
the individual}, and here the relation between the conscious and the
unconscious in us then appears in a peculiar light. The working of the
universal, although a spiritual working, becomes in relation to the
individual something partially or totally unconscious, both because the
universal is the comprehensive, the overarching, which is not
contained in the individual spirit, whose relation to it is the
relation of the single member of the totality to the totality itself
and its principle, --- and because the forms of activity of the
universal are not congruent with the forms of activity of finite
consciousness. Even when the individual spirit becomes most clearly
conscious of its relation to the pervading universal, the unconscious
working is converted into consciousness only within a certain limit,
because this very act is unconsciously steered and determined by what
works more deeply, just as the act by which the a priori was made an
object of consciousness was determined by the very working of the a
priori, unconscious to the individual.

The middle term through which the unconscious influence of the
principle of existence upon the individual life of consciousness is
realized \emph{is the real whole of existence}. This, as a totality
determined by the principle, is in its real being the rationally
determined, the planned (teleological); but it is the
\emph{unconsciously} rational, the \emph{unconsciously} thinking, the
\emph{unconsciously} purposive. The unconscious reason of nature is the
foundation of the conscious, individual life of reason. We think only
inasmuch as we have a brain that has been formed by the unconscious
thought of nature into
% --- p. 55 ---
\opage{55}the organ of our thought. Behind our conscious thought lies
``that which thinks in us''; and the involuntary, unconscious
production of thought proceeding from the essence, from the ``genius of
thought'', what we might call ``inner thinking'', ``thought in and for
itself'', lies behind the consciousness of this thinking, behind the
presentation of thought, not only to others but to ourselves.

In \emph{historical development} the reciprocal relation between the
conscious and the unconscious appears in pregnant forms. Since
historical development is a teleological development, and since in
such a development the goal of development is at every stage
\emph{ideally} present in its wholeness as that which determines and
guides the development, the steering and guiding power of development
is present as that which reaches beyond the individual consciousness,
and whose content is realized for consciousness only in a limited
form. In its wholeness it is therefore something unconsciously
working: it works behind consciousness, hidden from consciousness.
Therefore it works as the individual's \emph{fate}, which can indeed
be illuminated by consciousness, but only partially, since the
principle of thought becomes an object for thought according to a
\emph{single} determinateness, a \emph{single} basic relation, or only
in an \emph{abstract} way. The principle of historical development, in
its working as a concrete totality, preserves, as that which pervades
the consciousness of the individual, its stamp of a fate that is only
relatively explained as a conscious law of spirit. The universal
principle of spirit, at every moment, unconsciously for the
individual, gathers the individual consciousness together under its
overarching necessity.

\secnum{4.}

From one further standpoint the relation of conscious mental life to
unconscious mental life shows itself to us, when we see \emph{death}
as the return of the total life of consciousness into the unconscious.
Since our life of consciousness in all its forms \emph{is} throughout
organically conditioned, the dissolution of the organism must
necessarily entail the cessation of all these forms of consciousness,
and the nearest expression for this becomes: a return into the
unconscious.
% --- p. 56 ---
\opage{56}But this expression can conceal a double meaning. It can
designate a dissolution of the self into the material elements whose
forces worked together in the living individual. That would be
materialism's interpretation of the expression. But it can also
designate a return into a, as it were, more deepened and more
inwardized form of being, for which the determinations of the life of
consciousness no longer have validity. And we are directed back to
such an interpretation by the fact that the freedom of the human
spirit includes the guarantee of its eternity, secures it against
being able to vanish as something without essence; but that the same
freedom points back to the divine as its source, to self-surrender to
the principle, to union of essence with it as the highest goal, and
precisely thereby points beyond the sphere in which the life of
consciousness, with its relativity and finite oppositions, has a
place. But when we then ask about the form of the being of spiritual
life after the life of consciousness has vanished along with its
conditions, we raise a question which, precisely because in our
thinking we are bound to the forms of the life of consciousness, can
be answered by us only through hints and analogies. We then seek an
analogy in a relation that was touched on in the foregoing: in the
relation of the individual act of consciousness to the total
unconscious mental life and to the whole of the life of consciousness.
When the individual content of consciousness that has been called
forth by an external influence, a sense-image, for example, steps away
from the horizon of consciousness and goes back into the unconscious,
it enters through the latter into a comprehensive connection of
memories and thoughts; it is appropriated by this connection,
transformed, inwardized from sense-image into representation and
thought, and as such belongs to our comprehensive spiritual unity of
life. In this inwardization through the return into the unconscious we
might seek an analogy that can suggest the form of the transformation
of the individual self at the cessation of the life of consciousness.
Its going back into the unconscious would then be one with a
transformation from the form of phenomenal, sensibly determined
externality to a deeper inwardness. It would then
% --- p. 57 ---
\opage{57}be a moment in the life of the eternal, in the unity from
which the unfolding of life in existence eternally proceeds; and it
would be joined together with the other individual spirits in the
organization of the eternal fullness of life. That analogy can also be
extended to the relation between individual spirits. Just as our
individual acts of consciousness, in their conscious separation from
one another, have a continuity \emph{beneath the boundary line of
consciousness} and are therefore gathered in the comprehensive,
uniting unity of spirit, so the individual spirits, whose organisms are
indeed members of one comprehensive physical totality, can be
conceived as joined together in their unconscious ground by one that
bears, embraces and connects them, and as going back, [in]
% TR: the Danish line „Døden gaaende tilbage ...“ opens with a blank about one short word wide (a lost sort, evidently „i“); rendered in that sense, the missing word supplied as "[in]"
death, to this ground, and in it to a deeper, more inward union.

\begin{center}\rule{2cm}{0.4pt}\end{center}

Thus the life of consciousness would then become like a region in our
psychic life that was bounded on both sides by the unconscious. It
rises up out of the ground of unconscious \emph{organic} living; it is,
during the development of human life, at every moment borne and
conditioned by the unconscious, in constant reciprocal action with it;
and as it reaches its close, it goes back into an unconscious that now
expresses a higher \emph{spiritual} form of life, raised above the
limitation of the life of consciousness, above its separation of
externality and its finitude: a life in unity with the divine, which
in its eternal knowledge of itself is not, like the conscious self,
limited by something alien, but in this knowledge includes the whole
of existence, and includes it not as what is merely reflected, but as
that which has being only through the essence that in its knowledge
embraces it.

\begin{flushright}\textbf{H. Bröchner.}\end{flushright}
\begin{center}\rule{3cm}{0.4pt}\\[-10pt]\rule{3cm}{0.4pt}\end{center}

\end{document}
